\Needspace{34\baselineskip}
\begin{center}
GAUSS AN VON ZACH.
\end{center}

\begin{center}
Monatliche Correspondenz. Band XV. S. 152–157. Februar 1807.
\end{center}

\begin{flushright}
Braunschweig, 3. Januar 1807.
\end{flushright}

\par\indent --- --- --- Ich habe die Ehre, Ihnen hier die Ephemeride für den Lauf der \emph{Ceres} zu übersenden
die Elemente für ihre Berechnung konnte ich mangels anderer\HistoricalNote{26} nur auf die Beobachtung der letzten Opposition durch Herrn Professor
\SourceLeaf{286}%
\textsc{Pasquich} (M.\,C., Band~XIII, \href{https://www.digitale-sammlungen.de/en/view/bsb10538605?page=200}{S.~192}; Band~XIV, \href{https://www.digitale-sammlungen.de/en/view/bsb10538606?page=102}{S.~92}) stützen. Ihre weitere Verbesserung muss ich mir daher vorbehalten, bis ich zuverlässigere Beobachtungen erhalte.
Es waren die folgenden:

\begin{center}
XI. \emph{Elemente.}
\end{center}

\begin{center}\begingroup
\fontsize{9}{11}\selectfont
\setlength{\tabcolsep}{3pt}\renewcommand{\arraystretch}{1.12}
\begin{tabular}{ll}
\multicolumn{2}{l}{Epoche der Länge im Meridian von} \\
\quad Seeberg 1806 .\ .\ .\ .\ .\ .\ .\ . & $108^\circ\,19'\,34''7$ \\
Länge des Aphels 1806 .\ .\ . & $326\;37\;59$ \\
Länge des aufsteigenden Knotens$\cdot$1806  & $80\;53\;23$ \\
Neigung der Bahn .\ .\ .\ .\ .\ .\ . & $10\;37\;33.7$ \\
Mittlere jährliche tropische Bewegung  & $78^\circ\,9'\,23''3$ \\
Mittlere tägliche tropische Bewegung . & $770''8584$ \\
Exzentrizität .\ .\ .\ .\ .\ .\ .\ .\ .\ . & $0.0783486$ \\
Log. der Halbachse .\ .\ .\ .\ .\ . & $0.4420728$ \\
\end{tabular}\endgroup\end{center}

\par\indent Mein Werkchen über die Bestimmung der Planetenbahnen usw., dem ich den größten Teil meiner Zeit gewidmet habe, nähert sich der Vollendung. Bisher habe ich 22 Bogen Manuskript ausgearbeitet und bin gerade dabei, die ausführlichen Beispiele einzutragen. Eines stammt von der \emph{Juno}, bei der die Zwischenzeit 22 Tage beträgt. Beim zweiten, von der \emph{Pallas}, sind es 71 Tage. Das dritte, von der \emph{Ceres}, bei dem es 118 Tage waren, habe ich wieder gestrichen, um mit einem anderen, bei dem es 260 Tage sind, die große Allgemeinheit der Methode noch besser zu zeigen. In all diesen Beispielen lässt sie sich unmittelbar anwenden, ohne dass man irgendetwas über die Bahn als bereits bekannt voraussetzt. Das Ganze wird im Druck etwa 30 Bogen umfassen. Bei den gegenwärtigen Verhältnissen, fürchte ich, wird es in Deutschland vielleicht einige Mühe kosten, einen zuverlässigen Verleger zu finden. --- --- --- [Der geozentrische Lauf der Ceres vom 19. Dezember 1806 bis zum 21. September 1807 ist ebenso dargestellt, wie ich oben auf \hyperlink{native-283}{S.~278} dieses Bandes für die Pallas angedeutet habe. Nur ist die Lichtstärke auf fünf Delimalen genau bestimmt, dort waren es vier.]
